\section{Discussion}\label{sec:discussion}

\subsection{Why the Wavelet-Domain Cues Help}\label{subsec:fusion_insight}\label{subsec:freq_complement}\label{subsec:wavelet_choice}

The revised evidence supports a focused interpretation of the main result. Under the unified re-training protocol, WaveletMask improves the RTN on the synthetic benchmark (Table~\ref{tab:ablation_core}), so the most stable conclusion is not that every metric moves in the same direction, but that wavelet-domain priors improve the overall trade-off between structural repair and perceptual plausibility, which supports the proposed mask as an auxiliary restoration cue rather than a dataset-specific optimization artifact. Consistent with this, the DAVIS experiment (Table~\ref{tab:davis_cross_content}) changes the underlying scenes while keeping the degradation pipeline fixed, and the gain over the RTN persists, indicating that the cue is not tied to one video category, camera motion pattern, or clip identity, although cross-degradation generalization remains open.

\textbf{Why parameter-free maximum fusion is preferred.}
The fusion ablation in Table~\ref{tab:fusion_comparison} is best stated directly: among the tested strategies, the parameter-free branch-level maximum rule is the most reliable choice for the current learned-fusion design, the $b=+5$ spatial-prior initialization, the 200K-iteration training budget, and the reconstruction-only supervision, which is an architecture- and training-specific finding rather than a general verdict against learned fusion. The likely reason is the heterogeneity of the two inputs: LHFM is sparse and highly localized, whereas GFM is broader and tied to low-frequency photometric change, so a learned convex mixture can attenuate a sparse but informative response when the other input is weak but spatially widespread. The branch-level maximum rule avoids this mismatch by preserving the stronger cue at each location, which is especially useful when the two branches encode different temporal assumptions rather than two noisy estimates of the same quantity.

The same picture also clarifies why the flow-compensated (FC) variant does not surpass the uncompensated setting: compensating before mask formation can blur the inter-frame deviations that LHFM exploits and may inject optical-flow errors around occlusion boundaries, and because RTN-style propagation already performs temporal alignment downstream~\cite{rtn2022,rrtn2024}, the uncompensated mask preserves old-film defect contrast better, consistent with the motion-binned analysis in which the improvement is largest in the high-motion and high-occlusion terciles (Table~\ref{tab:motion_binning}).

\textbf{What the frequency-domain cues add.}
The frequency-domain design is most interpretable when LHFM and GFM are read as modeling two different classes of temporal irregularity. LHFM targets thin, spatially localized, high-frequency defects such as dust and narrow scratches: Haar detail sub-bands retain these abrupt local changes while suppressing much of the low-frequency scene content, so wavelet-domain inter-frame differences become more sensitive to defect-like perturbations than raw pixel differences, which matches the LHFM gains in Table~\ref{tab:ablation_core} and recent \textit{Applied Sciences} findings that frequency-domain representations are particularly effective for structured degradations~\cite{kim2025hybrid,park2025spatialfrequency}. GFM instead encodes gradual low-frequency photometric instability; its broader response profile is expected because flicker is a frame-level or region-level brightness inconsistency rather than a point artifact, which also explains why GFM tends to help SSIM more than PSNR in Table~\ref{tab:ablation_core}.

Two diagnostics reinforce this reading. The complementarity probe (Table~\ref{tab:cue_complementarity}) shows that the original RTN spatial indicator, LHFM, and GFM do not activate on identical regions and that the two wavelet branches are nearly disjoint in their top-response areas, precisely the regime in which branch-wise preservation is more appropriate than premature averaging. The residual-proxy probe (Table~\ref{tab:mask_interpretability}) shows strong alignment of GFM and the max-fused mask with the low-frequency proxy, while the weaker alignment of LHFM indicates that it functions as a gating prior for recurrent restoration rather than as a surrogate for supervised scratch or dust segmentation. This also positions WaveletMask relative to learned all-in-one restoration models such as AirNet, PromptIR, and MambaOFR~\cite{li2022airnet,potlapalli2023promptir,mao2025mambaofr}: instead of learning degradation representations at global, prompt, or state-space levels, WaveletMask injects a dense, hand-structured temporal cue into an RTN-style recurrent gate, trading generality for low overhead and interpretability.

\textbf{Wavelet family choice.}
Finally, the wavelet-family sweep (Table~\ref{tab:followup_ablations}) indicates that the benefit does not come from adding an arbitrary transform. Single-level Haar remains the most stable default under the current training budget, two-level Haar improves several perceptual metrics at a small PSNR cost, more oscillatory families such as db2 and sym2 are less stable in this lightweight setting, and DTCWT remains competitive but does not surpass Haar; for sparse old-film defects, sharp spatial localization is more valuable than a richer but less localized basis, and Haar's fixed $2\times2$, parameter-free filters match this goal. The cross-domain diagnostic (Table~\ref{tab:cross_domain_frequency_ablation}) narrows the claim: Fourier temporal responses and proxy-supervised shallow filters can align more strongly with residual-proxy labels, so Haar is the most stable \emph{tested} wavelet for end-to-end restoration gating rather than a universally optimal frequency detector. This is consistent with wavelet-based restoration work showing that decompositions can enlarge effective receptive fields~\cite{liu2018mwcnn} and that complex-valued transforms offer richer directional selectivity and approximate shift invariance~\cite{selesnick2005dual}, while for mask construction in old film restoration the simpler Haar basis already captures the defect statistics that matter most under the current recipe.

\subsection{Transferability, Boundary Conditions, and Evaluation Scope}\label{subsec:transferability}

Two challenging cases are best understood as boundary conditions rather than contradictions of the method. First, persistent fixed scratches are difficult for LHFM by construction. The detector compares short-term differences and is most sensitive to transient, non-stationary defects; if a scratch remains in nearly the same image-plane position for three or more frames, the geometric-mean difference is attenuated and the response is strongly reduced (Table~\ref{tab:persistent_scratch_probe}). Drifting scratches are less problematic because they continue to produce inter-frame inconsistency. The relevant conclusion is therefore not that LHFM fails on scratches in general, but that it faces challenges when the degradation becomes temporally stationary.

Second, GFM assumes gradual photometric evolution. Scene cuts and abrupt exposure flashes violate this assumption, so elevated GFM responses around such events are expected (Table~\ref{tab:gfm_failure_probe}). The luminance-difference guard already shows that most of this false activation can be suppressed without changing the backbone. Scene cuts should therefore be treated as a boundary condition of the current cue definition, not as evidence against low-frequency flicker modeling itself.

The no-reference evaluation on real footage should be interpreted with the same care. BRISQUE and NIQE are both natural-scene-statistics measures, but they are built on different priors: BRISQUE is opinion-aware and operates in the spatial domain~\cite{mittal2012brisque}, whereas NIQE is opinion-unaware and measures deviation from pristine-image statistics without distortion-specific training~\cite{mittal2013niqe}. CLIPIQA+, in turn, derives its preference from CLIP-based visual-language priors and is therefore more sensitive to semantic and perceptual plausibility~\cite{wang2023clipiqa}. Old-film restoration often changes grain, contrast, local sharpness, and semantic realism at the same time, so disagreement among BRISQUE, NIQE, and CLIPIQA+ is not only possible but expected. Recent \textit{Applied Sciences} work on no-reference assessment likewise emphasizes the limited universality of individual NR metrics and the value of combining complementary indicators instead of relying on a single score~\cite{rubel2022combinednr}. For this reason, the no-reference metrics are interpreted jointly with the qualitative comparisons rather than individually.

The scope of the comparison protocol should also be stated explicitly. All ten methods, including WaveletMask, were re-trained and re-evaluated under one unified benchmark pipeline (Table~\ref{tab:literature_context}), which supports a direct cross-method ranking on the synthetic split. To isolate the proposed gate from training confounds, RTN and WaveletMask were additionally compared under a strictly identical paired protocol (Appendix~\ref{app:stats}, Table~\ref{tab:same_protocol_main}). The two views are complementary: the controlled paired comparison remains the primary evidence for the proposed gate, while the unified benchmark situates the magnitude of the gains within the old film restoration landscape.

\subsection{Limitations and Future Work}\label{subsec:limitations}

Despite the encouraging evidence, several limitations define the current scope. First, all full-reference experiments rely on the SRWOV degradation pipeline. The DAVIS study confirms cross-content transfer, but transfer to independently synthesized degradations or real paired archival benchmarks remains open. Second, the present multi-seed analysis is stronger than a single-run report yet still modest; a larger five-seed or ten-seed study would provide tighter uncertainty estimates. Third, the mask is computed without explicit scene-cut awareness and without explicit mask supervision, so it may fail to detect temporally stationary damage and may produce false activation around abrupt editing events; in addition, clip-boundary frames lack bilateral temporal neighbors and are assigned zero wavelet masks, so the first and last frames of each clip currently receive no wavelet-domain degradation cue. Fourth, while the motion- and occlusion-binned analysis in Table~\ref{tab:motion_binning} shows that ordinary motion does not erode the restoration gain on the present benchmark, that evidence is camera-motion dominated and tercile-grained; object-dominant motion, extreme displacements, and absolute motion thresholds remain to be characterized. Fifth, the cross-domain diagnostic in Table~\ref{tab:cross_domain_frequency_ablation} shows that Fourier-temporal and proxy-supervised operators can align more strongly with residual proxies than frozen Haar responses, and these alternatives have not yet been retrained end-to-end inside the gating pipeline; the present evidence therefore supports Haar as the most stable tested wavelet choice rather than the optimal frequency operator. Finally, very large missing regions or long-term severe deterioration prove difficult for the current restoration-only formulation, because the problem gradually shifts from defect suppression to content completion.

These limitations suggest several concrete next steps. A first direction is event-aware gating, including scene-cut and flash detectors for GFM and persistence-aware detectors for stationary scratches. A second direction is to revisit learned fusion with stronger supervision, uncertainty calibration, and longer schedules, since the current max rule likely benefits from its robustness to scale mismatch more than from greater expressive power. A third direction is to couple restoration with explicit completion modules when film loss becomes spatially extensive; recent \textit{Applied Sciences} work on unified inpainting and colorization provides a useful template for this scenario~\cite{yeo2025integration}. Finally, broader validation on real archival collections and a larger public release of per-run evidence would strengthen both the empirical claims and the reproducibility value of the work.

Overall, the results support a balanced claim: WaveletMask is a lightweight and interpretable improvement for RTN-style old film restoration when degradations remain temporally informative, while persistent structural damage, abrupt temporal discontinuities, and large missing regions define the main frontier for future development.
